\begin{table}[htbp]
    \centering
    \caption{Matriks Keputusan Pemilihan Sistem Manajemen Basis Data}
    \label{tab:matriks_database}
    \begin{tabularx}{\textwidth}{|l|X|X|}
        \hline
        \textbf{Solusi} & \textbf{Kelebihan} & \textbf{Kekurangan} \\
        \hline
        PostgreSQL & Menjamin konsistensi data yang tinggi dan mendukung kueri relasional yang kompleks. Memiliki fitur analitik serta mekanisme \textit{indexing} yang sangat mumpuni untuk pengelolaan basis data skala besar. & Prosedur konfigurasi dan optimasi awal relatif lebih kompleks dibandingkan dengan MySQL. \\
        \hline
        MySQL & Memiliki tingkat kemudahan penggunaan yang baik, performa yang stabil untuk aplikasi umum, serta dukungan komunitas yang sangat luas. & Kapabilitas analitik dan fitur kueri tingkat lanjut (\textit{advanced query}) tidak sekomprehensif yang dimiliki oleh PostgreSQL. \\
        \hline
        MongoDB & Menawarkan fleksibilitas tinggi untuk skema data semi-terstruktur serta mendukung skalabilitas horizontal yang efisien melalui mekanisme \textit{sharding}. & Kurang optimal untuk menangani relasi data yang sangat kompleks serta membutuhkan manajemen khusus untuk menjamin konsistensi transaksi tingkat tinggi. \\
        \hline
    \end{tabularx}
\end{table}